% Results of the full-cycle note; \input by cycle_optimisation.tex.
% Numbers are macros from generated/macros.tex.

All results are for the LEI-V3 with the flight-calibrated (semi-empirical)
ROM, at the wind of its \code{cycle\_profile.yaml}: an \windLaw{} profile with
\SI{\windRef}{m/s} at \SI{\windHeight}{m}. They come from two different code
states, which this section keeps apart:
\begin{itemize}
  \item \textbf{Seeds} (Sect.~\ref{sec:res-seeds}) were generated on
        2026-10-06 with the ROM rework later merged in PR~\#22 (standardised
        steering, bound $|u_s|\le\limSteer$; AoA bound
        \SI{\limAoaMaxDeg}{\degree}), still on the semi-empirical ROM, i.e.\
        before the aerostructural ROM became the LEI-V3 default.
  \item \textbf{The optimum} (Sect.~\ref{sec:res-opt}) is the last converged
        LEI-V3 full-cycle solve, from 2026-08-27, at the same wind. It
        predates the Dubins reel-in, the reel-in design step and the steering
        standardisation, so it is shown as an illustration of what the staged
        NLP does, not as a result of the code described above.
\end{itemize}
Sect.~\ref{sec:res-status} states why no fresh optimum is reported.

\subsection{Seed generation}
\label{sec:res-seeds}

Two seeds with \dubinsHalves{} visible half-figures were generated. The
\emph{lobe} seed used \code{--auto}, starting from the default knobs of
Fig.~\ref{fig:lobe}. The \emph{Dubins} seed used \code{--reelin dubins
--loops~7 --check}: no tuning, only the forward simulation and the grid
recalibration. Figure~\ref{fig:res-seeds} shows both paths and
Table~\ref{tab:seeds} summarises them.

\begin{figure}[H]
  \centering
  \includegraphics{results_seeds.pdf}
  \caption{The two generated LEI-V3 seeds in the wind window; reel-in (depower
  above mid-band) in orange. Left: lobe seed after \code{--auto}; right:
  Dubins seed.}
  \label{fig:res-seeds}
\end{figure}

\begin{table}[H]
  \centering\small
  \caption{Generated seeds and their forward simulations. Steering and AoA
  ranges are to be read against the NLP bounds $|u_s|\le\limSteer$ (plus
  the \num{\steerSlack} warm-start slack) and
  $\alpha\le\SI{\limAoaMaxDeg}{\degree}$; closure against
  \SI{\closureTol}{m}.}
  \label{tab:seeds}
  % GENERATED by report_cycle_optimisation.py
\begin{tabular}{lrr}
\toprule
quantity & lobe seed & dubins seed \\
\midrule
$M$ / $N$ & $60$ / $283$ & $70$ / $275$ \\
nodes trimmed & $283/283$ & $275/275$ \\
max curvature (1/m) & $0.077$ & $0.065$ \\
cycle duration (s) & $94.2$ & $91.0$ \\
average power (kW) & $1.83$ & $1.78$ \\
closure $r_\mathrm{end}-r_0$ (m) & $+14.5$ & $+11.2$ \\
$u_s$ range (-) & $-0.22$ / $+0.19$ & $-0.19$ / $+0.19$ \\
max $\alpha$ (deg) & $7.4$ & $7.1$ \\
max tension (kN) & $3.96$ & $3.95$ \\
tuner iterations & 8 & -- \\
generation time (min) & $16$ & $5$ \\
\bottomrule
\end{tabular}

\end{table}

\paragraph{What worked.}
Both seeds trim at every node of the recalibrated grid, stay under the
curvature limit, and keep the angle of attack inside its bound with a wide
margin. The Dubins seed reached this \emph{without any tuning}. Its
steering grazes the bound (\seedDubinsSteerMin{}/\seedDubinsSteerMax) by less
than the slack. The lobe seed needed the tuner (Table~\ref{tab:trace}), and
its first three moves are the case the lever chains were written for. The
steering violation sat at the top of the reel-in ($s\approx0.01$--$0.02$; the
window is centred on $s=0$), and widening the bow $|a|$ shrank it step by
step, from 31 violating nodes to 9.

\paragraph{What did not.}
Neither seed closes to within \SI{\closureTol}{m}: they end
\SI{\seedLobeClosure}{m} (lobe) and \SI{\seedDubinsClosure}{m} (Dubins) long.
For the lobe seed, the tuner's closure lever was already exhausted at full
depower depth, and its fallback, shortening the reel-out fraction from
0.65 to 0.53, moved the gap by less than a metre. The tuner therefore
stopped and wrote its best iterate (``\autoLobeStuck''). A closure that is
insensitive to the reel-in fraction means that the reel-in speed, and not
the reel-in duration, limits how much tether comes back. The reel-in speed
is set by the depower-shifted winch law of Eq.~\eqref{eq:winch}. On paper, a
\SIrange{11}{17}{m} gap on $r_0=\SI{\seedRzero}{m}$ is a small closure
residual that the NLP's first iterations can absorb (Sect.~\ref{sec:seedreq}).

\begin{table}[H]
  \centering\small
  \caption{\code{--auto} trace of the lobe seed: one knob per iteration,
  chosen from the dominant symptom and its location.}
  \label{tab:trace}
  % GENERATED by report_cycle_optimisation.py
\begin{tabular}{rp{0.5\textwidth}l}
\toprule
it. & dominant symptom & knob move \\
\midrule
1 & steering [-0.37, +0.19] at 31 node(s), worst at s=0.02 & az\_reelin\_amp $-0.36\to-0.42$ \\
2 & steering [-0.32, +0.19] at 12 node(s), worst at s=0.01 & az\_reelin\_amp $-0.42\to-0.48$ \\
3 & steering [-0.24, +0.19] at 9 node(s), worst at s=0.01 & az\_reelin\_amp $-0.48\to-0.54$ \\
4 & under-reel +14.9 m at full depower & reelout\_fraction $0.65\to0.61$ \\
5 & under-reel +14.8 m at full depower & reelout\_fraction $0.61\to0.57$ \\
6 & steering [-0.26, +0.18] at 9 node(s), worst at s=0.02 & az\_reelin\_amp $-0.54\to-0.60$ \\
7 & under-reel +14.5 m at full depower & reelout\_fraction $0.57\to0.53$ \\
8 & cannot close the cycle (gap +15.8 m) within knob ranges & stop, best iterate \\
\bottomrule
\end{tabular}

\end{table}

\begin{figure}[H]
  \centering
  \includegraphics{results_seed_series.pdf}
  \caption{Forward simulations of the two seeds, which are the NLP warm
  starts. Dotted: NLP bounds on steering and angle of attack. Both cycles
  are about \SI{90}{s} long; the reel-in is the stretch of negative $v_r$ and
  high $u_p$.}
  \label{fig:res-seedseries}
\end{figure}

\subsection{An optimised cycle (2026-08-27)}
\label{sec:res-opt}

Figures~\ref{fig:res-archived-power} and~\ref{fig:res-archived} show the last
converged solve. The figure was written by that run itself and compares its
seed (a lobe seed of that time) with the optimum. The staged NLP raised the
cycle-average power from \SI{1.24}{kW} to \SI{\archPower}{kW} and shortened
the cycle to \SI{\archDuration}{s}, with the closure satisfied to
\SI{\archClosure}{m} and $r_0$ moved to \SI{\archRzero}{m}. The mechanisms
are visible in the per-node decisions:
\begin{itemize}
  \item \textbf{Reel-out at higher load.} The depower profile leaves the
        identified band towards the powered side ($u_p$ down to
        \SI{\archDepMin}{m}), and the tension rises to \SI{\archTensionMax}{kN}
        peak (\SI{\archTensionMean}{kN} mean), reeling out at up to
        \SI{\archVrMax}{m/s}.
  \item \textbf{Shorter, faster reel-in.} The reel speed reaches
        \SI{\archVrMin}{m/s}, and the reel-out time fraction of the cycle is
        \num{\archReeloutFraction}.
  \item \textbf{Free-speed winch.} The optimal $T(v_r)$ points do not lie
        on the flight-regressed force law the seed was marched with
        (Fig.~\ref{fig:res-archived-power}, bottom left). The regression of a
        controller from the optimum is a separate step.
\end{itemize}

\begin{figure}[H]
  \centering
  \includegraphics[width=\textwidth]{archived_power.png}
  \caption{Seed versus optimum of the 2026-08-27 LEI-V3 solve, as written by
  \code{run\_full\_cycle\_opti.py}: instantaneous and average power, cumulative
  energy, winch operating plane and metrics. Shaded: reel-in.}
  \label{fig:res-archived-power}
\end{figure}

\begin{figure}[H]
  \centering
  \includegraphics{results_archived.pdf}
  \caption{The 2026-08-27 optimum ($M=\archM$, $N=\archNodes$ nodes). Left:
  optimised spline and the node positions, reel-in ($v_r<0$) in orange.
  Right: reel speed, tension, depower and path rate $\dot s$; the dotted line
  is the numerical floor $\dot s\ge0.01$.}
  \label{fig:res-archived}
\end{figure}

\paragraph{Caveats of this optimum.}
Four features of the 2026-08-27 solution show what the constraints of that
time did not yet cover. Each motivated a later addition to the method:
\begin{enumerate}
  \item The depower profile goes down to \SI{\archDepMin}{m}, below the band
        $[\upPowered,\upDepowered]$\,m on which the ROM was identified, so
        the reel-out extrapolates the ROM. The band is now declared in the
        ROM file, and the AoA bound can be set from the ROM's validity range.
  \item The flown turn radius drops to \SI{\archTurnRadiusMin}{m}, because
        no turn-radius constraint existed. The optional sub-sampled
        constraint of Table~\ref{tab:constraints} was added for this.
  \item \archFloorNodes{} of the \archNodes{} nodes sit on the numerical
        floor $\dot s=0.01$: \archFloorReelin{} while reeling in, most of
        the rest in the transition into the reel-in (Fig.~\ref{fig:res-archived},
        bottom right). There the time step $\Delta t=\Delta s/\dot s$ is
        capped by a guard rather than set by the physics, so the reel-in
        timing of this optimum is partly an artefact of that bound.
  \item One node ($s=\archSdotMaxS$) has $\dot s=\SI{\archSdotMax}{s^{-1}}$,
        about 150 times the median \SI{\archSdotMedian}{s^{-1}}. Since
        $v_\tau=r\sigma'\dot s$ stays ordinary, the arc rate $\sigma'$ of the
        spline nearly vanishes there and $\dot s$ jumps to compensate, a
        \emph{hesitation point}. The floor $\sigma'\ge0.2\bar\sigma'$ that now
        comes with the turn-radius constraint
        (Table~\ref{tab:constraints}) excludes it.
\end{enumerate}
Its steering also reaches $|u_s|=\archSteerMax$, which was the steering bound
before the steering input was standardised.

\subsection{Status}
\label{sec:res-status}

With the code of 2026-10-06, both on commit \code{722488a} and with the ROM
rework of PR~\#22, the staged NLP does not converge for the LEI-V3 at this
wind. Started from the seeds above, the main stage stays close to feasible
(primal infeasibility \num{0.1}--\num{0.4}) for about 200 iterations, then
diverges and ends in restoration at the \maxIter-iteration cap. On the last
commit the seed itself is harder: the tuner lowered $\beta_\mathrm{pk}$ to
its floor without bringing the reel-in angle of attack inside its bound. Where
the regression since 2026-08-27 comes from (ROM, steering standardisation,
seed changes or staging) has not been isolated. Until it has, the seed
results of Sect.~\ref{sec:res-seeds} describe the current code, and the
optimum of Sect.~\ref{sec:res-opt} describes what the method produced when
it last converged.
